\scanpage{47}
\mspage{23}
homomorphismes de groupes profinis associés, \deleted{regarder} \inserted{étudier} la situation analogue pour des ``modèles'' élémentaires\inserted{ anabéliens}, voire des modèles $K(\pi,1)$ généraux. (On peut aussi regarder des tels modèles sur un corps \deleted{\unsure{de nombres}} \inserted{$K$ ext.} de type fini de $\mathbb{Q}$ -- mais passons pour le moment sur cette situation mixte, un peu bâtarde\ldots) Si $U,V$ sont des tels modèles, tout morphisme $\unsure{V|U}\longrightarrow\unsure{U|V}$ définit un morphisme de topos galoisiens sur $B_{\mathbb{Q}}$, $B_U\longrightarrow B_V$, \deleted{\unsure{de tels}} et si $V$ est \deleted{anabélien} élémentaire anabélien, cet homomorphisme est connu quand on connaît seulement
\[
H_1(B_{\bar U},\mathbb{Z}_{\ell})-H_{\unsure{1}}(B_{\bar V},\mathbb{Z}_{\ell})
\]
-- ce qui est beaucoup moins que la classe d'iso. du hom. de $B_{\mathbb{Q}}$-topos. (En fait, \deleted{il suffit} \inserted{\unsure{sans hyp.}} anabélienne sur $V$, dès que $V$ \unsure{se plonge|est plongé} dans une VA, $f$ est connu quand on connaît son action sur les \unsure{topos|types} étales\ldots). Mais quels sont les hom. $B_U\longrightarrow B_V$, ou $E_U\longrightarrow E_V$
\margin{46}
